% GENERATED FILE. Edit canonical lessons, then run npm run generate:books.
% canonical-source-hash: fnv1a64:556447d5034b6331
% canonical-lessons: AR-C95-jayb, AR-C95-khatam, AR-C95-qubbaa, AR-C95-hidha, AR-C95-mitaf

\chapter{Pocket, Ring, Hat, Shoe, Coat}
\label{ch:ar-pocket-ring-hat-shoe-coat}
\hlchaptermodality{95}

\begin{chapteropening}
\textbf{By the end of this chapter:} \emph{I can say a pocket, a ring, a hat, a shoe and a coat.}

Complete the last lesson of chapter 95: I can say a pocket, a ring, a hat, a shoe and a coat.
\end{chapteropening}

\section[\texorpdfstring{\ar{جيب} (jayb)}{جيب (jayb)}]{\ar{جيب} (jayb) — a pocket}

\label{lesson:AR-C95-jayb}



\begin{quote}
Before the new one: say the Arabic for trousers, then the Arabic for a sock.
\end{quote}

\subsection*{You'll want to know: \ar{جيب}}
\textbf{\ar{جيب}} — \emph{jayb} — \textquotedblleft{}a pocket\textquotedblright{}.

\begin{cousinweb}
\textbf{\ar{جيب}} (\emph{jayb}) is built on the native Arabic root \emph{j-y-b}, \textquotedblleft{}the opening of a shirt at the neck\textquotedblright{} (Lane). Lane's \emph{jayb} is the opening at the neck and bosom of a shirt; \textquotedblleft{}pocket\textquotedblright{} came later. Medieval Latin translators read the Arabic word for a sine, a loan from Sanskrit \emph{jyā}, as this word and rendered it \emph{sinus}, the source of English \emph{sine}.
\end{cousinweb}

\begin{grammarlens}[title={What you've built}]
One more word for this chapter.
\end{grammarlens}

\subsection*{Your turn}
\begin{itemize}
\raggedright
  \item \emph{Say it:} \emph{jayb}
  \item \emph{Say it:} \emph{jayb}, once more
  \item \emph{Say it:} say \emph{jayb}
  \item \emph{Recall:} say the Arabic for trousers, then the Arabic for a sock, then say \emph{jayb} again
\end{itemize}

\begin{tcolorbox}[breakable,colback=teal!4,colframe=teal!35!black,title={Before you move on}]
What does \ar{جيب} mean? (\textquotedblleft{}A pocket\textquotedblright{}.) Say it once more.
\end{tcolorbox}
\section[\texorpdfstring{\ar{خاتم} (khātam)}{خاتم (khātam)}]{\ar{خاتم} (khātam) — a ring}

\label{lesson:AR-C95-khatam}



\begin{quote}
Before the new one: say the Arabic for a sock, then the Arabic for a pocket.
\end{quote}

\subsection*{You'll want to know: \ar{خاتم}}
\textbf{\ar{خاتم}} — \emph{khātam} — \textquotedblleft{}a ring\textquotedblright{}.

\begin{cousinweb}
\textbf{\ar{خاتم}} (\emph{khātam}) is built on the native Arabic root \emph{kh-t-m}, \textquotedblleft{}to seal\textquotedblright{} (Lane). \emph{Khātam} is a seal, a signet ring.
\end{cousinweb}

\begin{grammarlens}[title={What you've built}]
One more word for this chapter.
\end{grammarlens}

\subsection*{Your turn}
\begin{itemize}
\raggedright
  \item \emph{Say it:} \emph{khātam}
  \item \emph{Say it:} \emph{khātam}, once more
  \item \emph{Say it:} say \emph{khātam}
  \item \emph{Recall:} say the Arabic for a sock, then the Arabic for a pocket, then say \emph{khātam} again
\end{itemize}

\begin{tcolorbox}[breakable,colback=teal!4,colframe=teal!35!black,title={Before you move on}]
What does \ar{خاتم} mean? (\textquotedblleft{}A ring\textquotedblright{}.) Say it once more.
\end{tcolorbox}
\section[\texorpdfstring{\ar{قبعة} (qubbaʿa)}{قبعة (qubbaʿa)}]{\ar{قبعة} (qubbaʿa) — a hat}

\label{lesson:AR-C95-qubbaa}



\begin{quote}
Before the new one: say the Arabic for a pocket, then the Arabic for a ring.
\end{quote}

\subsection*{You'll want to know: \ar{قبعة}}
\textbf{\ar{قبعة}} — \emph{qubbaʿa} — \textquotedblleft{}a hat\textquotedblright{}.

\begin{cousinweb}
\textbf{\ar{قبعة}} (\emph{qubbaʿa}) is usually traced to Syriac \emph{qubbaʿā}, \textquotedblleft{}cap\textquotedblright{} (EtymDB); Lane has no entry for it, though a root \emph{q-b-ʿ} exists in his lexicon.
\end{cousinweb}

\begin{grammarlens}[title={What you've built}]
One more word for this chapter.
\end{grammarlens}

\subsection*{Your turn}
\begin{itemize}
\raggedright
  \item \emph{Say it:} \emph{qubbaʿa}
  \item \emph{Say it:} \emph{qubbaʿa}, once more
  \item \emph{Say it:} say \emph{qubbaʿa}
  \item \emph{Recall:} say the Arabic for a pocket, then the Arabic for a ring, then say \emph{qubbaʿa} again
\end{itemize}

\begin{tcolorbox}[breakable,colback=teal!4,colframe=teal!35!black,title={Before you move on}]
What does \ar{قبعة} mean? (\textquotedblleft{}A hat\textquotedblright{}.) Say it once more.
\end{tcolorbox}
\section[\texorpdfstring{\ar{حذاء} (ḥidhāʾ)}{حذاء (ḥidhāʾ)}]{\ar{حذاء} (ḥidhāʾ) — a shoe}

\label{lesson:AR-C95-hidha}



\begin{quote}
Before the new one: say the Arabic for a ring, then the Arabic for a hat.
\end{quote}

\subsection*{You'll want to know: \ar{حذاء}}
\textbf{\ar{حذاء}} — \emph{ḥidhāʾ} — \textquotedblleft{}a shoe\textquotedblright{}.

\begin{cousinweb}
\textbf{\ar{حذاء}} (\emph{ḥidhāʾ}) is built on the native Arabic root \emph{ḥ-dh-w}, \textquotedblleft{}to cut a sandal to a pattern\textquotedblright{} (Lane). Lane's verb \emph{ḥadhā} is to measure and cut a sandal by a pattern.
\end{cousinweb}

\begin{grammarlens}[title={What you've built}]
One more word for this chapter.
\end{grammarlens}

\subsection*{Your turn}
\begin{itemize}
\raggedright
  \item \emph{Say it:} \emph{ḥidhāʾ}
  \item \emph{Say it:} \emph{ḥidhāʾ}, once more
  \item \emph{Say it:} say \emph{ḥidhāʾ}
  \item \emph{Recall:} say the Arabic for a ring, then the Arabic for a hat, then say \emph{ḥidhāʾ} again
\end{itemize}

\begin{tcolorbox}[breakable,colback=teal!4,colframe=teal!35!black,title={Before you move on}]
What does \ar{حذاء} mean? (\textquotedblleft{}A shoe\textquotedblright{}.) Say it once more.
\end{tcolorbox}
\section[\texorpdfstring{\ar{معطف} (miʿṭaf)}{معطف (miʿṭaf)}]{\ar{معطف} (miʿṭaf) — a coat}

\label{lesson:AR-C95-mitaf}



\begin{quote}
Before the new one: say the Arabic for a hat, then the Arabic for a shoe.
\end{quote}

\subsection*{You'll want to know: \ar{معطف}}
\textbf{\ar{معطف}} — \emph{miʿṭaf} — \textquotedblleft{}a coat\textquotedblright{}.

\begin{cousinweb}
\textbf{\ar{معطف}} (\emph{miʿṭaf}) is built on the native Arabic root \emph{ʿ-ṭ-f}, \textquotedblleft{}to bend, to incline\textquotedblright{} (Lane). Lane's \emph{miʿṭaf} is a cloak, the same as \emph{ʿiṭāf}.
\end{cousinweb}

\begin{grammarlens}[title={What you've built}]
One more word for this chapter.
\end{grammarlens}

\subsection*{Your turn}
\begin{itemize}
\raggedright
  \item \emph{Say it:} \emph{miʿṭaf}
  \item \emph{Say it:} \emph{miʿṭaf}, once more
  \item \emph{Say it:} say \emph{miʿṭaf}
  \item \emph{Recall:} say the Arabic for a hat, then the Arabic for a shoe, then say \emph{miʿṭaf} again
\end{itemize}

\begin{tcolorbox}[breakable,colback=teal!4,colframe=teal!35!black,title={Before you move on}]
What does \ar{معطف} mean? (\textquotedblleft{}A coat\textquotedblright{}.) Say it once more.
\end{tcolorbox}
